
\subsection{A spectral cap from spin flip}
\label{sec:fqe-spectral-spinflip}

The following bound controls the part of a cut spectrum above $1/2$
directly by two-copy quantities. It holds for every pure state, not only on
the critical locus, and needs neither a four-qubit normal form nor a
polynomial certificate.
Write $x_+=\max\{x,0\}$ and use the Hilbert--Schmidt norm throughout.

\begin{proposition}[Spectral spin-flip bound]
\label{prop:fqe-spectral-spinflip}
For every unit pure four-qubit state and every balanced cut $A|A^c$,
\begin{equation}\label{eq:fqe-spectral-spinflip}
 \boxed{\quad
 (2\lambda_{\max}(\rho_A)-1)_+^2
 \le \tau+\frac12W_1
 =2(P-1)-\frac12W_1\le2(P-1).
 \quad}
\end{equation}
Both coefficients on the first right-hand side are optimal. In particular,
$P=1$ forces every balanced eigenvalue to be at most $1/2$.
\end{proposition}

The matrix inequality behind the proposition holds in every bipartite
dimension. The two vectors below need only have the same Schmidt
coefficients; their Schmidt bases may be unrelated.

\begin{lemma}[Overlap and marginal mismatch]
\label{lem:fqe-isospectral-overlap}
Let $\psi,\phi$ be unit bipartite pure states with common squared Schmidt
coefficients $\lambda_1\ge\cdots\ge\lambda_n\ge0$. Let
$\rho_A,\rho_B$ and $\sigma_A,\sigma_B$ be their marginals. Then
\begin{equation}\label{eq:fqe-isospectral-overlap}
 |\langle\psi,\phi\rangle|^2+
 \frac12\bigl(\|\rho_A-\sigma_A\|_F^2+
                    \|\rho_B-\sigma_B\|_F^2\bigr)
 \ge (2\lambda_1-1)_+^2.
\end{equation}
For each common Schmidt spectrum with $\lambda_1\ge1/2$, the bound is
attained by changing the sign of every Schmidt term except the first.
\end{lemma}

\begin{proof}
The case $\lambda_1\le1/2$ is immediate. Assume $\lambda_1>1/2$.
Pad the coefficient matrices with zeros if the local dimensions differ,
and choose Schmidt bases for $\psi$. Its coefficient matrix is
$D=\operatorname{diag}(\sqrt{\lambda_i})$; write $W$ for that of $\phi$.
Thus $WW^\dagger$ and $W^\dagger W$ have eigenvalues $\lambda_i$.
Put
\[
 \Lambda=D^2,\qquad r=1-\lambda_1,\qquad d=2\lambda_1-1,
 \qquad S=\operatorname{diag}(0,r-\lambda_2,\ldots,r-\lambda_n).
\]
The entries of $S$ are nonnegative and increasing. Trace rearrangement
therefore gives, for $R=WW^\dagger$ and for $R=W^\dagger W$,
\begin{equation}\label{eq:fqe-isospectral-rearrange}
 \operatorname{tr}(SR)\ge
 \sum_{i=2}^n(r-\lambda_i)\lambda_i
 =r^2-\sum_{i=2}^n\lambda_i^2.
\end{equation}
Indeed, diagonalising $R$ expresses its diagonal as a doubly stochastic
average of permutations of the decreasing vector $(\lambda_i)$; pairing
it with the increasing diagonal of $S$ is minimised in the displayed order.

Define the quadratic form
\[
 \mathcal B_D(W)=\|W\|_F^2+|\operatorname{tr}(DW)|^2
       -\operatorname{tr}\bigl(\Lambda(WW^\dagger+W^\dagger W)\bigr).
\]
Its useful lower bound has an explicit positive remainder. Let
$W_{\rm t}=(W_{ij})_{i,j\ge2}$ and
$t=\sum_{i=2}^n\sqrt{\lambda_i}W_{ii}$, so $\operatorname{tr}(DW)=\sqrt{\lambda_1}W_{11}+t$. Since
$\Lambda+S=rI+dE_{11}$, with $E_{11}$ the first matrix unit, and $1-2r=d$, the left side below is
$d\bigl(\|W_{\rm t}\|_F^2-|W_{11}|^2\bigr)+|\sqrt{\lambda_1}W_{11}+t|^2$; for $r>0$, completing the square in
$W_{11}$ with $\lambda_1-d=r$ gives
\begin{equation}\label{eq:fqe-isospectral-calibration}
 \begin{split}
 \mathcal B_D(W)-\operatorname{tr}\bigl(S(WW^\dagger+W^\dagger W)\bigr)
 ={}&d\left(\|W_{\rm t}\|_F^2-\frac{|t|^2}{r}\right)
       +\frac{|rW_{11}+\sqrt{\lambda_1}t|^2}{r}\ge0.
 \end{split}
\end{equation}
Cauchy--Schwarz gives
$|t|^2\le r\sum_{i\ge2}|W_{ii}|^2\le r\|W_{\rm t}\|_F^2$.
This proves positivity without any condition on $W$.
When $r=0$, the remainder is simply $\|W_{\rm t}\|_F^2\ge0$.

Apply~\eqref{eq:fqe-isospectral-rearrange} twice to obtain
\[
 \mathcal B_D(W)\ge
 2\left(r^2-\sum_{i=2}^n\lambda_i^2\right).
\]
Set $c_2=\sum_i\lambda_i^2$. In these coordinates
$\langle\psi,\phi\rangle=\operatorname{tr}(DW)$, $\rho_A=\Lambda$, $\sigma_A=WW^\dagger$,
and $\rho_B,\sigma_B$ are the transposes of $\Lambda,W^\dagger W$. Since
$\operatorname{tr}(WW^\dagger)^2=\operatorname{tr}(W^\dagger W)^2=c_2$, the mismatch terms
are $2c_2-2\operatorname{tr}(\Lambda WW^\dagger)$ and
$2c_2-2\operatorname{tr}(\Lambda W^\dagger W)$. With $\|W\|_F^2=1$, the left-hand
side of~\eqref{eq:fqe-isospectral-overlap} is therefore
$\mathcal B_D(W)-1+2c_2$. Its lower bound is
$2r^2+2\lambda_1^2-1=d^2$, as required. Taking
$W=\operatorname{diag}(\sqrt{\lambda_1},-\sqrt{\lambda_2},\ldots,
-\sqrt{\lambda_n})$ proves sharpness.
\end{proof}

\begin{proof}[Proof of Proposition~\ref{prop:fqe-spectral-spinflip}]
Apply Lemma~\ref{lem:fqe-isospectral-overlap} to
$\phi=\widetilde\psi=\sigma_y^{\otimes4}\overline\psi$.
Spin flip preserves every Schmidt spectrum and
$|\langle\psi,\widetilde\psi\rangle|^2=\tau$.
For a two-qubit marginal on sites $i,j$, Pauli expansion gives
\[
 \rho_{ij}-\widetilde\rho_{ij}
 =\frac12\left(\mathbf r_i\cdot\boldsymbol\sigma\otimes I
              +I\otimes\mathbf r_j\cdot\boldsymbol\sigma\right),
 \qquad
 \|\rho_{ij}-\widetilde\rho_{ij}\|_F^2
 =\|\mathbf r_i\|^2+\|\mathbf r_j\|^2.
\]
Thus the two marginal mismatch terms add to $W_1/2$.
The purity identity gives the remaining equalities and inequality
in~\eqref{eq:fqe-spectral-spinflip}.

For sharpness take
$\psi=(a|00\rangle+b|11\rangle)_{12}\otimes
(|00\rangle+|11\rangle)_{34}/\sqrt2$, with real $a,b\ge0$ and
$a^2+b^2=1$. On $12|34$, the left side is one, while
$\tau=4a^2b^2$ and $W_1=2(a^2-b^2)^2$, so equality holds throughout.
The choices $a=b$ and $b=0$ force respectively the coefficients of
$\tau$ and $W_1$ to be at least $1$ and $1/2$.
\end{proof}

Equation~\eqref{eq:fqe-spectral-spinflip} supplies a fully analytic quantum
compatibility constraint. The separate-cut bounds do not by themselves
replace the quadratic variance certificate: SCB-5 gives exact classical spectra
satisfying all three caps and even $F\ge\tau$ (hence the critical
envelope~\eqref{eq:fqe-critical-envelope}) but violating (QK), so compatibility
between the three cut spectra still matters.


\subsection{Quantitative rigidity of two maximally mixed cuts}
\label{sec:fqe-renyi-two-flat}

\begin{proposition}[Crossed purity costs and two-flat rigidity]
\label{prop:fqe-two-flat-rigidity}
Let $\psi$ be any unit pure four-qubit state. Fix a balanced cut $k$,
write its ordered spectrum as
$\lambda_1\ge\lambda_2\ge\lambda_3\ge\lambda_4\ge0$,
and put
\[
\Delta_k=\sum_{j\ne k}(p_j-1/4),\qquad
p_j=\operatorname{tr}\rho_{1j}^2.
\]
If $\lambda_1\ge1/2$, then
\[
\boxed{\quad
\Delta_k\ge\frac14W_1+
2(\lambda_2\lambda_3+\lambda_2\lambda_4+
\lambda_3\lambda_4)
\ge\frac14W_1+(\lambda_3+\lambda_4)^2.\quad}
\tag{TF}\label{eq:fqe-two-flat}
\]
The first inequality is sharp for every chosen spectrum with
$\lambda_1\ge1/2$. On the critical locus, equality holds exactly when
all nonzero squares of the other normal-form amplitudes have a phase
opposite to the square of the largest amplitude. The noncritical
sharpness family in Proposition~\ref{prop:fqe-spectral-spinflip} also
attains the first inequality.

Consequently, whenever \emph{two} balanced cuts of a pure
four-qubit state have maximally mixed marginals $(1/4)^4$,
the third has Schmidt rank at most two. Up to local
unitaries and a permutation of the qubits, every such
state belongs to the attained cluster-to-Bell family
\eqref{eq:fqe-quadratic-family}. In particular the exact
R\'enyi envelope of Proposition~\ref{prop:fqe-renyi-rank-two} of the supplement
applies for all $q>0$.
\end{proposition}
\begin{proof}
The purity identity and Proposition~\ref{prop:fqe-spectral-spinflip} give
\[
2\Delta_k
=2(P-p_k-1/2)
=W_1+\tau-(2p_k-1)
\ge\frac12W_1+(2\lambda_1-1)^2-(2p_k-1).
\]
Using $\sum_i\lambda_i=1$ gives
\[
(2\lambda_1-1)^2-(2p_k-1)
=4(\lambda_2\lambda_3+\lambda_2\lambda_4+
\lambda_3\lambda_4).
\]
This proves the first bound. Since
$\lambda_2\ge(\lambda_3+\lambda_4)/2$,
the second follows by dropping the nonnegative
$\lambda_3\lambda_4$ term.
On the critical locus, write the chosen spectrum as $\lambda_i=|w_i|^2$
in the normal form, so $\tau=|\sum_iw_i^2|^2$.
The bound is then the reverse triangle inequality, with exactly the stated
opposite-phase equality condition. Choosing those phases realises every
ordered spectrum with $\lambda_1\ge1/2$, proving sharpness.

Suppose now two cuts have purity $1/4$. Already one
maximally mixed two-qubit cut forces all four one-qubit
marginals to be $I/2$ by purity of complementary
reductions. If $k$ is the remaining cut, then
$p_k=P-1/2=1/2+\tau/2\ge1/2$. Since
$p_k=\sum_i\lambda_i^2\le\lambda_1$, the hypothesis
$\lambda_1\ge1/2$ holds. The first bound and
$\Delta_k=0$ force at most one among
$\lambda_2,\lambda_3,\lambda_4$ to be nonzero.
The opposite-phase equality condition means that, up to
overall phase and a signed permutation, $w$ has the form
\[
w=(\sqrt a,\ i\sqrt{1-a},\ 0,\ 0),\qquad 1/2\le a\le1.
\]
The magic-basis action is
$\operatorname{diag}(w)\mapsto O_1\operatorname{diag}(w)O_2^T$, $O_1,O_2\in SO(4)$: the same permutation on both
sides permutes the entries of $w$, a sign on one side changes the sign of one entry, and extra signs on the zero
coordinates make both factors proper rotations. This action is induced by local one-qubit unitaries on
each side of the chosen cut, so this is exactly the
normal form \eqref{eq:fqe-quadratic-family} up to LU and
cut relabelling.
\end{proof}

